\documentclass[11pt]{article}
\usepackage[margin=1in]{geometry}
\usepackage{amsmath,amssymb,amsthm}
\usepackage{array}
\usepackage{parskip}
\usepackage{booktabs}

\newtheorem{theorem}{Theorem}
\newtheorem{lemma}[theorem]{Lemma}
\newtheorem{corollary}[theorem]{Corollary}
\newtheorem{proposition}[theorem]{Proposition}
\theoremstyle{definition}
\newtheorem{definition}[theorem]{Definition}
\newtheorem{example}[theorem]{Example}
\theoremstyle{remark}
\newtheorem{remark}[theorem]{Remark}

\newcommand{\F}{\mathbb{F}}
\newcommand{\PG}{\mathrm{PG}}
\newcommand{\AG}{\mathrm{AG}}

\title{Disproof of conjecture \texttt{00000008434}}
\author{earthking11}
\date{2026-09-14}

\begin{document}
\maketitle

\section*{Verdict: FALSE}

We disprove conjecture \texttt{00000008434} as stated. The first clause is
refuted by an explicit, genuine block design: $\AG(3,2)$, the unique
$3$-$(8,4,1)$ Steiner quadruple system, has intersection spectrum
$\{0,2\}$, whereas the interval $[s_0,s_1] = [0,2]$ also contains the value
$1$, which never occurs. The second clause is refuted by the Pasch
configuration, a $1$-$(6,3,2)$ design with $s_0 = s_1 = 1$ that is not
symmetric. Both refutations are finite, elementary, and completely explicit;
the first is additionally verified inside Lean~4.

\section{The conjecture, quoted verbatim}

From \texttt{conjectures/00000008434.md}:

\begin{quote}
\textbf{English.} Definition: The intersection structure of a combinatorial
geometry: the intersection spectrum of its blocks. Conjecture: The
intersection spectrum is exactly the set of all integer values in the
interval $[s_0,s_1]$; and $s_0 = s_1$ if and only if the design is symmetric.
(interval characterization of intersection spectra)
\end{quote}

The file is bilingual. The Chinese text is reproduced verbatim (in Unicode) in
\texttt{README.md}; the two versions have identical content. The PDF font does
not embed CJK glyphs, so the Chinese is kept in the markdown copy rather than
duplicated here, and nothing is lost by working from the English below.

\section{The statement defines none of its terms}

The quoted statement is not a formal mathematical sentence: it introduces
names (``combinatorial geometry'', ``blocks'', ``intersection spectrum'',
``symmetric'') without defining any of them, and it never says what $s_0$ and
$s_1$ are. A disproof therefore has to \emph{pin a reading}, and we record the
one adopted here. It is the only standard reading under which the wording has
content:

\begin{definition}[Pinned reading]\label{def:pinned}
A \emph{combinatorial geometry} is a block design: a finite point set $P$ with
a family $\mathcal{B}$ of \emph{blocks} (a set system on $P$). Its
\emph{intersection spectrum} is the set of sizes of intersections of
\emph{distinct} blocks,
\[
S \;=\; \{\, |B \cap B'| \;:\; B,B' \in \mathcal{B},\ B \neq B' \,\},
\]
and $s_0 := \min S$, $s_1 := \max S$. A design is \emph{symmetric} when the
number of blocks equals the number of points, $b = |\mathcal{B}| = |P| = v$
(equivalently, for $2$-designs, $r = k$).
\end{definition}

Two points about this reading deserve emphasis.

\begin{enumerate}
\item Taking $s_0,s_1$ to be the \emph{minimum and maximum of the spectrum} is
  forced by the phrasing: the conjecture asserts an ``interval
  characterization'' whose endpoints are $s_0,s_1$, and the second clause
  ``$s_0 = s_1$ iff symmetric'' only has content when $s_0,s_1$ are the two
  ends of $S$. (Under this reading that second clause is in fact a
  \emph{theorem} for genuine $2$-designs; see
  Proposition~\ref{prop:secondclause}. It is the generality to
  $1$-designs that kills it.)
\item The refutation of clause~1 does not depend on any looseness in the
  reading. The witness is a full-fledged $3$-$(8,4,1)$ design --- the strongest
  design hypothesis available in this context --- so no ``intended''
  strengthening of ``combinatorial geometry'' can rescue the claim.
\end{enumerate}

\section{The witness: the geometry $\AG(3,2)$}

\begin{definition}[The $14$ affine planes]\label{def:ag32}
Let $P = \F_2^3$, the $8$ points. For a nonzero normal vector
$a \in \F_2^3 \setminus \{0\}$ and $b \in \{0,1\}$ put
\[
B_{a,b} \;=\; \{\, x \in \F_2^3 \;:\; a \cdot x = b \,\},
\qquad
a\cdot x = a_1x_1 + a_2x_2 + a_3x_3 \pmod 2 .
\]
The point $x=(x_1,x_2,x_3)$ is encoded by the $8$-bit mask with bit
$x_1 + 2x_2 + 4x_3$ set, so that bitwise \texttt{AND} computes intersections.
The family $\mathcal{B} = \{B_{a,b} : a \neq 0,\ b \in \{0,1\}\}$ consists of
$7 \cdot 2 = 14$ blocks; in sorted mask form they are
\[
[\,15, 51, 60, 85, 90, 102, 105, 150, 153, 165, 170, 195, 204, 240\,].
\]
\end{definition}

\begin{proposition}\label{prop:sqs}
$\mathcal{B}$ is a genuine $3$-$(8,4,1)$ design: $v = 8$, $b = 14$,
$k = 4$, every $3$-subset of the points lies in exactly one block.
Consequently each point lies in $r = 7$ blocks and each pair in
$\lambda_2 = (v-2)/(k-2) = 3$ blocks, and $\lambda_3 = 1$.
\end{proposition}

\begin{proof}
For fixed $a \neq 0$, the map $x \mapsto a\cdot x$ is a nonzero $\F_2$-linear
functional, so each $B_{a,b}$ is a coset of the kernel $\ker(a\cdot)$, a
$2$-dimensional subspace; hence $|B_{a,b}| = 2^2 = 4$ and the $14$ blocks are
affine $2$-planes.

Let $T \subset P$ with $|T| = 3$. A line of $\AG(3,2)$ has exactly two points,
so no three distinct points are collinear: the three points are affinely
independent and span a unique affine $2$-subspace, i.e.\ a unique set of the
form $B_{a,b}$. Hence exactly one block contains $T$, and
$\lambda_3 = 1$. The stated $\lambda_2,\lambda_1$ follow from the standard
identities
\[
\lambda_2 = \frac{v-2}{k-2}\,\lambda_3 = \frac{6}{2} = 3,
\qquad
\lambda_1 = \frac{v-1}{k-1}\,\lambda_2 = \frac{7}{3}\cdot 3 = 7 .
\]
Alternatively, both were checked exhaustively by \texttt{reproduce.py} and in
Lean (all $56$ triples, all $28$ pairs, all $8$ points).
\end{proof}

This is the unique $3$-$(8,4,1)$ design (the Steiner quadruple system
$\mathrm{SQS}(8)$, equivalently the affine geometry $\AG(3,2)$); the
uniqueness is classical and is not needed here.

\section{The intersection spectrum of $\AG(3,2)$}

\begin{proposition}\label{prop:spectrum}
For two distinct blocks $B_{a,b} \neq B_{a',b'}$ of Definition~\ref{def:ag32}:
\[
|B_{a,b} \cap B_{a',b'}|
=
\begin{cases}
0, & a = a' \ (\text{the two planes are parallel}),\\[2pt]
2, & a \neq a' \ (\text{distinct normals}).
\end{cases}
\]
Hence the spectrum is $S = \{0,2\}$, with $s_0 = \min S = 0$ and
$s_1 = \max S = 2$.
\end{proposition}

\begin{proof}
If $a = a'$ then, since the blocks are distinct, $b \neq b'$; the equations
$a\cdot x = b$ and $a\cdot x = b'$ are incompatible, so
$B_{a,0} \cap B_{a,1} = \varnothing$. There are exactly $7$ nonzero normals
$a$, and each gives exactly one such disjoint pair, so this case occurs
$7$ times.

If $a \neq a'$, they are linearly independent, so the two affine planes have
distinct directions; their intersection is a nonempty affine subspace of
dimension $2+2-3 = 1$, i.e.\ an affine line, which in $\AG(3,2)$ has exactly
$2$ points. This case occurs for all the remaining pairs, of which there are
\[
\binom{14}{2} - 7 = 91 - 7 = 84 .
\]
Thus the spectrum, taken over all $91$ unordered pairs of distinct blocks,
contains only $0$ and $2$; both values occur, so $S=\{0,2\}$.
\end{proof}

The seven disjoint pairs are exactly the seven parallel classes
$\{B_{a,0},B_{a,1}\}$, $a \neq 0$; in mask terms:
\[
(15,240),\ (51,204),\ (60,195),\ (85,170),\ (90,165),\ (102,153),\ (105,150).
\]

\subsection*{The $91$-pair spectrum table}

The next table is the complete intersection-size data. The left block lists,
for each candidate intersection size $m \in \{0,1,2\}$, the number of unordered
pairs $\{B,B'\}$ of distinct blocks with $|B\cap B'| = m$; the right block
records whether $m$ lies in the conjectured interval $[s_0,s_1]=[0,2]$.

\begin{center}
\begin{tabular}{r r c}
\toprule
intersection size $m$ & number of pairs with $|B\cap B'| = m$ & $m \in [s_0,s_1]=[0,2]$? \\
\midrule
$0$ & $7$  & yes \\
$1$ & $0$  & yes \\
$2$ & $84$ & yes \\
\midrule
total & $91 = \binom{14}{2}$ & \\
\bottomrule
\end{tabular}
\end{center}

So exactly $7 + 0 + 84 = 91$ pairs were examined, and the spectrum is
$S = \{0,2\}$: the value $1$ is \emph{absent}. The full $14 \times 14$
intersection-size matrix (upper triangle; diagonal blank) is displayed below,
so that every one of the $91$ pairs can be read off directly. Its entries are
$2$ everywhere off the diagonal except in the seven positions
$(1,14),(2,13),(3,12),(4,11),(5,10),(6,9),(7,8)$, which are $0$ --- the seven
parallel classes.

\begin{center}
\footnotesize
\setlength{\tabcolsep}{4pt}
\begin{tabular}{r|*{14}{c}}
\multicolumn{1}{c}{} & $1$ & $2$ & $3$ & $4$ & $5$ & $6$ & $7$ & $8$ & $9$ & $10$ & $11$ & $12$ & $13$ & $14$ \\
\hline
$1$  & & 2 & 2 & 2 & 2 & 2 & 2 & 2 & 2 & 2 & 2 & 2 & 2 & 0 \\
$2$  & & & 2 & 2 & 2 & 2 & 2 & 2 & 2 & 2 & 2 & 2 & 0 & 2 \\
$3$  & & & & 2 & 2 & 2 & 2 & 2 & 2 & 2 & 2 & 0 & 2 & 2 \\
$4$  & & & & & 2 & 2 & 2 & 2 & 2 & 2 & 0 & 2 & 2 & 2 \\
$5$  & & & & & & 2 & 2 & 2 & 2 & 0 & 2 & 2 & 2 & 2 \\
$6$  & & & & & & & 2 & 2 & 0 & 2 & 2 & 2 & 2 & 2 \\
$7$  & & & & & & & & 0 & 2 & 2 & 2 & 2 & 2 & 2 \\
$8$  & & & & & & & & & 2 & 2 & 2 & 2 & 2 & 2 \\
$9$  & & & & & & & & & & 2 & 2 & 2 & 2 & 2 \\
$10$ & & & & & & & & & & & 2 & 2 & 2 & 2 \\
$11$ & & & & & & & & & & & & 2 & 2 & 2 \\
$12$ & & & & & & & & & & & & & 2 & 2 \\
$13$ & & & & & & & & & & & & & & 2 \\
$14$ & & & & & & & & & & & & & & \\
\end{tabular}
\end{center}

\section{Clause 1 is false}

\begin{theorem}\label{thm:clause1}
The intersection spectrum $S$ of $\AG(3,2)$ is \emph{not} the set of all
integer values in the interval $[s_0,s_1]$. Indeed
\[
S = \{0,2\}
\qquad\text{but}\qquad
[s_0,s_1] \cap \mathbb{Z} = \{0,1,2\},
\]
so the value $1$ lies in the conjectured interval and is missing from the
spectrum.
\end{theorem}

\begin{proof}
Immediate from Proposition~\ref{prop:spectrum}: $s_0 = 0$, $s_1 = 2$, so the
integer interval is $\{0,1,2\}$, while the spectrum omits $1$; $S$ is a
\emph{strict} subset of the interval, with $|S| = 2 < 3 = |[s_0,s_1]\cap
\mathbb{Z}|$.
\end{proof}

\begin{corollary}
Conjecture \texttt{00000008434}, first clause, is false.
\end{corollary}

\begin{remark}
The failure is as small as it can be: the interval has length $2$ and exactly
one interior value, and that one value is the missing one. The witness is not
a boundary case of the definitions --- it is the classical
$3$-$(8,4,1)$ design, and the failure comes from clean geometry: two planes in
$\AG(3,2)$ are either parallel (intersecting in $0$ points) or meet in an
affine line ($2$ points); there is no way for them to meet in exactly
$1$ point.
\end{remark}

\section{Clause 2: false for $1$-designs, but a theorem for $2$-designs}

The second clause of the conjecture is ``$s_0 = s_1$ if and only if the
design is symmetric''. It holds for genuine $2$-designs, but the conjecture
says ``design'', and in the general $1$-design setting it fails.

\begin{example}[Pasch]\label{ex:pasch}
Let $P = \{1,2,3,4,5,6\}$ and
\[
\mathcal{B} = \bigl\{\{1,2,3\},\ \{1,4,5\},\ \{2,4,6\},\ \{3,5,6\}\bigr\}.
\]
This is the Pasch configuration. It has $b = 4$ blocks of size $k = 3$ on
$v = 6$ points; every point lies in exactly $r = 2$ blocks, so it is a
$1$-$(6,3,2)$ design. Every two distinct blocks meet in exactly one point,
\[
\{1,2,3\}\cap\{1,4,5\}=\{1\},\quad
\{1,2,3\}\cap\{2,4,6\}=\{2\},\quad
\{1,2,3\}\cap\{3,5,6\}=\{3\},
\]
and similarly for the other three pairs; hence $s_0 = s_1 = 1$. But
$b = 4 \neq 6 = v$, so the design is \emph{not} symmetric.
\end{example}

\begin{corollary}\label{cor:clause2}
Clause 2 fails in the generality stated: $s_0 = s_1$ does not imply that the
design is symmetric.
\end{corollary}

\begin{proposition}[The $2$-design case, stated to be fair]\label{prop:secondclause}
Let $D$ be a nontrivial $2$-$(v,k,\lambda)$ design (so $b \geq v$ by Fisher's
inequality, and the incidence matrix $N$ has rank $v$). If every two distinct
blocks of $D$ meet in the same number $\mu$ of points, then $D$ is symmetric,
i.e.\ $b = v$.
\end{proposition}

\begin{proof}
Let $N$ be the $b \times v$ incidence matrix. Since the off-diagonal entries of
$NN^{\mathsf T}$ are the pairwise intersection sizes, the hypothesis says
$NN^{\mathsf T} = (k-\mu) I_b + \mu J_b$. This matrix has eigenvalues
$k-\mu$ (multiplicity $b-1$) and $k-\mu+\mu b = k + \mu(b-1) > 0$. Since
$\operatorname{rank}(NN^{\mathsf T}) = \operatorname{rank}(N) = v < b$ when
$b > v$, the matrix would have to be singular, forcing $k - \mu = 0$, i.e.\
$\mu = k$: any two distinct blocks would share all $k$ points of a block,
hence be equal, contradicting $b > 1$. Therefore $b > v$ is impossible, and
with Fisher's inequality $b \geq v$ we get $b = v$.
\end{proof}

Thus the second clause is genuinely true for the symmetric-design setting it
was presumably intended for; it is the $1$-design generality of the wording
that makes the ``if and only if'' false. We report this to delimit the
refutation honestly: clause 2 is a defect of the statement's generality, while
clause 1 is refuted even in the strongest, most well-behaved setting.

\section{Reproducibility}

The complete finite verification is carried out by \texttt{reproduce.py}
using the Python~3 standard library only:

\begin{quote}\ttfamily\raggedright
\$ python3 reproduce.py\\
... \textit{(builds the 14 planes of $\F_2^3$, verifies the $3$-$(8,4,1)$\\
property, computes all 91 pairwise intersections and their multiplicities,\\
checks the Pasch $1$-design)}\\
PASS: all checks verified; conjecture 00000008434 is FALSE.
\end{quote}

It builds the planes from the affine equations (not from the hard-coded mask
list), asserts that the generated masks equal the pinned list, verifies all
$56$ triples have containment count $1$, tabulates the $91$-pair spectrum,
asserts $S = \{0,2\}$ and that $1 \in [0,2] \setminus S$, and verifies the
Pasch failure. It prints \texttt{PASS} and exits $0$ exactly when every check
holds (a buggy check was found and fixed this way during preparation).

The LaTeX document is rebuilt with

\begin{quote}\ttfamily\raggedright
\$ tectonic --outdir build main.tex
\end{quote}

producing \texttt{build/main.pdf}.

\section{Formalisation}

The refutation of clause 1 is formalised in the Lean~4 project under
\texttt{lean4/} (core Lean only, \texttt{import Std}, no Mathlib, no
\texttt{sorry}). Blocks are encoded as $8$-bit \texttt{Nat} masks, the
intersection as \texttt{Nat.land}, and the spectrum by a structurally
recursive function (so it reduces in the kernel). The main statements are

\begin{quote}\ttfamily\raggedright
theorem blocks\_length : blocks.length = 14\\
theorem triples\_length : triples.length = 56\\
theorem each\_triple\_in\_one\_block :\\
\hspace*{2em}triples.all (fun T => countIn T == 1) = true\\
theorem spectrum\_length : (spec blocks).length = 91\\
theorem spectrum\_values :\\
\hspace*{2em}(spec blocks).all (fun s => s == 0 || s == 2) = true\\
theorem spec\_count\_zero : ... filter (s == 0) ... length = 7\\
theorem spec\_count\_one  : ... filter (s == 1) ... length = 0\\
theorem spec\_count\_two  : ... filter (s == 2) ... length = 84\\
theorem s0\_eq : s0 = 0\\
theorem s1\_eq : s1 = 2\\
theorem intervalFull\_false : intervalFull = false\\
theorem intervalSpectrum\_false : \(\neg\) IntervalSpectrum\\
theorem second\_clause\_fails : ... \(\wedge\) \(\neg\) Symmetric pasch.length 6
\end{quote}

All are proved by kernel \texttt{decide} (with
\texttt{set\_option maxRecDepth 1000000}); there is no \texttt{sorry}.
\texttt{lean4/Check.lean} prints \texttt{\#print axioms} for every theorem; the
audit reports only \texttt{propext} (and no axioms at all for the pure
\texttt{Nat} equalities), and in particular no \texttt{sorryAx}. See
\texttt{lean4/README.md} for the statement table and the scope note.

\section{Caveats}

\begin{itemize}
\item \textbf{The file defines nothing.} Conjecture \texttt{00000008434} is a
  prose sentence: it fixes no precise meaning for ``combinatorial geometry'',
  ``blocks'', ``intersection spectrum'', ``symmetric'', and it never defines
  $s_0$ or $s_1$. We pinned the standard reading (Definition~\ref{def:pinned});
  the refutation of clause 1 is robust under any reasonable variant because
  the witness is a genuine $3$-design.
\item \textbf{Clause 2 is only false for $1$-designs.} For genuine $2$-designs
  it is a theorem (Proposition~\ref{prop:secondclause}). We flag this rather
  than overstate the refutation.
\item \textbf{Uniqueness of the witness.} $\AG(3,2)$ is the unique
  $3$-$(8,4,1)$ design; that classical fact is not required. Everything used is
  established directly by the finite computations.
\item \textbf{Lean scope.} The Lean development certifies the finite
  combinatorial facts (design property, spectrum, absence of $1$, negation of
  clause 1, and the Pasch failure). The general $2$-design theorem
  (Proposition~\ref{prop:secondclause}) is proved on paper in this document and
  is not formalised; it is not needed for the disproof.
\end{itemize}

\section*{Status against the submission rules}

Rule 3 requires a LaTeX source, a PDF document, and a Lean~4 project.

\begin{itemize}
\item \textbf{LaTeX source} --- \texttt{main.tex}, a standalone
  \texttt{article} using \texttt{geometry}, \texttt{amsmath},
  \texttt{amssymb}, \texttt{amsthm}, \texttt{array}, \texttt{parskip},
  \texttt{booktabs}; compiles with
  \texttt{tectonic --outdir build main.tex}.
\item \textbf{PDF} --- at \texttt{build/main.pdf}, produced by the command
  above.
\item \textbf{Lean~4 project} --- \texttt{lean4/} with \texttt{lean-toolchain}
  pinned to \texttt{leanprover/lean4:v4.33.1}, \texttt{lakefile.toml}
  (name \texttt{tlmc8434}, library \texttt{Main}, no dependencies),
  \texttt{Main.lean}, \texttt{Check.lean} (\texttt{\#print axioms}), and
  \texttt{lean4/README.md}. \texttt{lake build} exits $0$ and the audit
  reports no \texttt{sorryAx}.
\end{itemize}

\end{document}
